\section{Experiments}

\subsection{Validation Strategy: Synthetic + PoC Approach}

\textbf{Rationale:} Real Waterbirds experiments require GPU compatibility resolution (PyTorch 2.1+ for H100 sm\_90 architecture) or CPU training ($\sim$70-90 GPU hours). We validate methodology via:

\begin{enumerate}
\item \textbf{Synthetic validation (h-e1, h-m-integrated):} Proves pipeline correctness and mechanism plausibility using controlled synthetic gradients.
\item \textbf{Proof-of-concept (h-m-mitigate):} Demonstrates mitigation effectiveness on MNIST+Color toy dataset.
\end{enumerate}

\textbf{What synthetic validation proves:}
Code executes without errors (GradCAM $\rightarrow$ GAIA-Z $\rightarrow$ statistical testing); statistical tests compute correctly (divergence, p-values, effect sizes); mechanism holds when gradient differences exist (causality test works).

\textbf{What synthetic validation cannot prove:}
Real minority gradients exhibit abnormality (requires real gradient extraction); real spatial regularization improves WGA (requires real training with penalty).

\textbf{Epistemic status:} Methodology validated (HIGH confidence), real-world empirical claims pending (LOW confidence until real validation).

\subsection{Hypothesis Breakdown}

Three sub-hypotheses test different aspects of the gradient abnormality framework:

\textbf{h-e1 (Detection - MUST\_WORK gate):}
Gradient abnormality methodology differentiates minority vs majority patterns when synthetic gradients exhibit controlled near-zero rate differences.
\textbf{Success:} Divergence $\geq$0.2, $p<0.01$, Cohen's $d\geq$0.8
\textbf{Failure action:} Abandon gradient abnormality approach

\textbf{h-m-integrated (Mechanism - MUST\_WORK gate):}
Three experiments validate causal chain Steps 1-3:
\begin{itemize}
\item \textbf{Exp1:} GAIA divergence correlates with WGA ($|\rho|>0.7$, $p<0.05$)
\item \textbf{Exp2:} Background swap reduces GAIA-Z by $\geq$30\% (paired t-test)
\item \textbf{Exp3:} Minority accuracy $\geq$60\% (GradCAM validity check)
\end{itemize}
\textbf{Failure action:} Mechanism falsified, redesign regularization

\textbf{h-m-mitigate (Mitigation - SHOULD\_WORK gate):}
Spatial regularization improves WGA on MNIST+Color toy dataset by $\geq$10\% over ERM baseline.
\textbf{Success:} WGA improvement $\geq$10\%, no catastrophic accuracy drop
\textbf{Failure action:} PoC ineffective, defer mitigation to future work

\subsection{Experiment Design: h-e1 (Detection)}

\subsubsection{Synthetic Gradient Generation}

\textbf{Setup:} Mimic Waterbirds test set structure (5794 samples, 4 groups) with controlled GAIA-Z properties.

\textbf{Generation process:}
\begin{enumerate}
\item \textbf{Majority samples ($n=2900$ per group, $\sim$90\%):} Sample gradients from $\mathcal{N}(0, 0.5)$ with 60\% near-zero elements (simulates sparse gradients)
\item \textbf{Minority samples ($n=500$ per group, $\sim$10\%):} Sample gradients from $\mathcal{N}(0, 1.0)$ with 30\% near-zero elements (simulates dense gradients)
\item \textbf{Near-zero threshold:} $\epsilon = 10^{-5}$ (GAIA-Z standard)
\end{enumerate}

\textbf{Expected GAIA-Z scores:}
Majority: mean $\approx$ 0.60 (60\% near-zero rate); Minority: mean $\approx$ 0.30 (30\% near-zero rate); Divergence: 0.30 ($\geq$0.2 criterion)

\textbf{Validation:} Synthetic data proves pipeline \emph{can} differentiate patterns when differences exist. Does NOT prove real Waterbirds gradients exhibit this property.

\subsubsection{Metrics and Analysis}

\textbf{Primary metric:} GAIA-Z divergence = $|\text{mean(minority)} - \text{mean(majority)}|$
\textbf{Statistical test:} Welch's t-test ($H_0: \mu_{\text{minority}} = \mu_{\text{majority}}$)
\textbf{Effect size:} Cohen's $d$ (0.2=small, 0.5=medium, 0.8=large)

\subsection{Experiment Design: h-m-integrated (Mechanism)}

\subsubsection{Experiment 1: Correlation-WGA-GAIA Link}

\textbf{Setup:} Synthetic model sweep with varying WGA levels.

\textbf{Procedure:}
\begin{enumerate}
\item Generate 2 synthetic ``models'' with different WGA values (50\%, 90\%)
\item Assign GAIA divergence inversely proportional to WGA (higher divergence $\rightarrow$ lower WGA)
\item Compute Pearson correlation between divergence and WGA
\end{enumerate}

\textbf{Expected result:} $|\rho| > 0.7$, $p<0.05$ (strong correlation)

\subsubsection{Experiment 2: Background Augmentation Causality Test}

\textbf{Setup:} 100 synthetic minority samples with controlled GAIA-Z scores.

\textbf{Procedure:}
\begin{enumerate}
\item Assign original GAIA-Z scores (30\% near-zero rate)
\item Simulate background swap: increase near-zero rate to 60\% (mimics majority pattern)
\item Paired t-test: $H_0: \text{mean(original - augmented)} = 0$
\end{enumerate}

\textbf{Expected result:} Mean reduction $\geq$30\%, $p<0.05$, Cohen's $d\geq$0.5

\subsubsection{Experiment 3: Minority Accuracy Check}

\textbf{Setup:} Synthetic minority group with controlled accuracy.

\textbf{Procedure:}
\begin{enumerate}
\item Assign 68\% minority accuracy ($>$60\% threshold, mimics Waterbirds)
\item Verify GradCAM validity assumption (A1)
\end{enumerate}

\textbf{Expected result:} Accuracy $\geq$60\% (A1 satisfied)

\subsection{Experiment Design: h-m-mitigate (Mitigation PoC)}

\subsubsection{MNIST+Color Toy Dataset}

\textbf{Dataset construction:}
Base: MNIST digits (0-9); Spurious feature: Background color (red/blue); Correlation: 90\% (digit 0-4 on red, 5-9 on blue in training); Test: Balanced (50\% correlation, minority groups exist)

\textbf{Why MNIST?} Simple spurious feature (color) allows rapid PoC validation. Complex spurious (Waterbirds backgrounds) deferred.

\textbf{Subsample:} 10\% of MNIST (6000 train, 1000 test) for smoke test speed.

\subsubsection{Training Configuration}

\textbf{Baselines:}
\textbf{ERM:} Standard cross-entropy, no regularization;
\textbf{Spatial Regularization:} $L_{\text{CE}} + \lambda \cdot L_{\text{spatial}}$ (Section 3.4)

\begin{table}[h]
\centering
\caption{MNIST+Color training hyperparameters}
\label{tab:mnist_config}
\begin{tabular}{ll}
\toprule
Parameter & Value \\
\midrule
Architecture & 3-layer CNN (simple) \\
Optimizer & Adam \\
Learning rate & $10^{-3}$ \\
Batch size & 64 \\
Epochs & 2 (smoke test) \\
$\lambda_{\text{init}}$ & 0.01 \\
Percentile & 75 \\
Seed & 42 \\
\bottomrule
\end{tabular}
\end{table}

\textbf{Metrics:}
Worst-group accuracy (WGA): $\min(\text{accuracy across 4 groups})$;
Average accuracy: $\text{mean(accuracy across all samples)}$

\textbf{Success criterion:} WGA $\geq$ baseline + 10\%, average accuracy drop $\leq$2\%

\subsubsection{Limitations Acknowledged}

\textbf{L1 (Smoke test only):} 1 seed, 2 epochs (not 5 seeds $\times$ 50 epochs). No statistical significance testing (no bootstrap, no confidence intervals). PoC validates methodology, not robustness.

\textbf{L2 (Toy dataset):} MNIST color spurious simpler than Waterbirds backgrounds. Generalization to real complex spurious unknown.

\textbf{L3 (No GroupDRO comparison):} ERM baseline only. Competitive positioning requires GroupDRO/JTT comparison on Waterbirds.
